\chapter{更完整物理条件下的控制优化}

恢复重力和手部碰撞后重新验证接触方案，同时明确物理步长敏感性。动作块执行接口独立分析，不将解析轨迹当成学习策略。

\section{恢复重力与手部碰撞}

\label{stage:s13}

\subsection{反馈对应的改动}\label{s13-detail-1}

\FloatBarrier\begingroup\small\setlength{\tabcolsep}{2.5pt}
\begin{longtable}[]{@{}
  >{\raggedright\arraybackslash}p{(\linewidth - 4\tabcolsep) * \real{0.2299}}
  >{\raggedright\arraybackslash}p{(\linewidth - 4\tabcolsep) * \real{0.4373}}
  >{\raggedright\arraybackslash}p{(\linewidth - 4\tabcolsep) * \real{0.3327}}@{}}
\caption{反馈对应的改动：设置与结果}\label{tab:s13-1}\tabularnewline
\toprule\noalign{}
\begin{minipage}[b]{\linewidth}\raggedright
上次问题
\end{minipage} & \begin{minipage}[b]{\linewidth}\raggedright
该阶段处理
\end{minipage} & \begin{minipage}[b]{\linewidth}\raggedright
评价方式
\end{minipage} \\
\midrule\noalign{}
\endfirsthead
\toprule\noalign{}
\begin{minipage}[b]{\linewidth}\raggedright
上次问题
\end{minipage} & \begin{minipage}[b]{\linewidth}\raggedright
该阶段处理
\end{minipage} & \begin{minipage}[b]{\linewidth}\raggedright
评价方式
\end{minipage} \\
\midrule\noalign{}
\endhead
\bottomrule\noalign{}
\endlastfoot
无重力工况 & 启用重力；关节PID增加PhysX重力补偿向量；保留无补偿失败对照 & 前三轴5s误差、开始阶段接触情况 \\
只有球形探针 & 恢复panda\_hand与左右手指碰撞；求和三刚体对平面的接触力 & 核对场景、接触体和合力 \\
接触瞬间峰值偏大 & 接近减速、接触门控、冻结名义位置、平滑力目标、修正速度限制 & 1ms原始峰值、相对首次接触曲线 \\
对高度变化敏感 & 实际平面±1/±5mm扰动，减速器仍使用名义高度 & 同口径峰值和末段带内率 \\
材质过于理想 & 刚性与等效柔顺材质分组；摩擦系数变化 & 不把软材质结果当作刚性控制提升 \\
只满足短窗口 & 同时统计末2s均力带内率、原始力带内率和RMSE & 保留连续0.5s指标，不再只靠它下结论 \\
\end{longtable}\endgroup

\subsection{控制与接触模型}\label{s13-detail-2}

前5秒仍完成前三轴五次轨迹，增量{[}0.06,0.04,-0.04{]}rad。普通PID增益沿用上阶段配置，力矩限幅前四轴87Nm、后三轴12Nm，夹爪5N。主要条件控制周期10ms、物理1ms；后续单独改变控制周期，另有物理步长0.5ms的数值检查，评价仍按10ms分箱。

\begin{Verbatim}[breaklines=true,breakanywhere=true,fontsize=\footnotesize]
tau = Kp*(q_ref-q) - Kd*dq + Ki*integral + g(q)
M_adm * x_ddot + B_adm * x_dot = F_ref - F_feedback
z_ref = z_frozen - x
F_ref = 5 * (3*u**2 - 2*u**3),  u=clip((t-t_contact)/1s,0,1)
\end{Verbatim}

重力补偿使用本地Isaac 5.1关节接口返回的PhysX补偿向量，正号加入输出力矩。新增日志保存各关节补偿力矩，并检查调用前后关节速度未被改变。没有将真实Franka接口的隐含补偿假设直接套用进仿真。

接触发生后以实测手部高度冻结名义位置，沿法向叠加导纳位移。微分逆解正则化系数0.01，每次每轴参考增量限制0.003rad，另设总逆解跳变保护。三个接触刚体共用平面过滤器，力取世界Z分量之和。夹爪保持张开，不执行夹物动作；其他连杆碰撞尚未恢复。

减速器利用已知场景的名义平面高度和初始碰撞体边界估算间隙，接近速度随间隙降低，保留0.3mm/s下限。高度扰动不传给减速器。由于手指姿态会改变最低点，这个估计不是精确距离传感器，测试结果必须包含几何先验误差的影响。

初始对照采用16次位置求解迭代、4次速度迭代；最终统一方案改用4次位置迭代，并在相同设置下做扰动验证。迭代次数的影响单列讨论，不将这一变化归因于控制算法。球形探针与原装夹爪的接触几何不同，因此不将两组力误差的差值单独归因于算法。柔顺接触参数为等效刚度和阻尼，未经真实材料标定。

相关代码在 \path|scripts/run.py|；实验入口为 \path|scripts/suite.py|、\path|scripts/validation_suite.py| 和 \path|scripts/refine_suite.py|。开源设计对照见 \path|OPEN_SOURCE_REVIEW.md|，主要参考Franka官方力控制示例和NVIDIA稳定性、接触材质文档；未增加ROS控制器依赖。

\subsection{重力及几何对照}\label{s13-detail-3}

\FloatBarrier\begingroup\small\setlength{\tabcolsep}{2.5pt}
\begin{longtable}[]{@{}
  >{\raggedright\arraybackslash}p{(\linewidth - 12\tabcolsep) * \real{0.1783}}
  >{\raggedright\arraybackslash}p{(\linewidth - 12\tabcolsep) * \real{0.1348}}
  >{\raggedright\arraybackslash}p{(\linewidth - 12\tabcolsep) * \real{0.1348}}
  >{\raggedright\arraybackslash}p{(\linewidth - 12\tabcolsep) * \real{0.1348}}
  >{\raggedright\arraybackslash}p{(\linewidth - 12\tabcolsep) * \real{0.1348}}
  >{\raggedright\arraybackslash}p{(\linewidth - 12\tabcolsep) * \real{0.1348}}
  >{\raggedright\arraybackslash}p{(\linewidth - 12\tabcolsep) * \real{0.1477}}@{}}
\caption{重力及几何对照：设置与结果}\label{tab:s13-2}\tabularnewline
\toprule\noalign{}
\begin{minipage}[b]{\linewidth}\raggedright
条件
\end{minipage} & \begin{minipage}[b]{\linewidth}\raggedright
原始峰值/N
\end{minipage} & \begin{minipage}[b]{\linewidth}\raggedright
末段均力/N
\end{minipage} & \begin{minipage}[b]{\linewidth}\raggedright
力RMSE/N
\end{minipage} & \begin{minipage}[b]{\linewidth}\raggedright
周期带内率
\end{minipage} & \begin{minipage}[b]{\linewidth}\raggedright
原始带内率
\end{minipage} & \begin{minipage}[b]{\linewidth}\raggedright
四项检查
\end{minipage} \\
\midrule\noalign{}
\endfirsthead
\toprule\noalign{}
\begin{minipage}[b]{\linewidth}\raggedright
条件
\end{minipage} & \begin{minipage}[b]{\linewidth}\raggedright
原始峰值/N
\end{minipage} & \begin{minipage}[b]{\linewidth}\raggedright
末段均力/N
\end{minipage} & \begin{minipage}[b]{\linewidth}\raggedright
力RMSE/N
\end{minipage} & \begin{minipage}[b]{\linewidth}\raggedright
周期带内率
\end{minipage} & \begin{minipage}[b]{\linewidth}\raggedright
原始带内率
\end{minipage} & \begin{minipage}[b]{\linewidth}\raggedright
四项检查
\end{minipage} \\
\midrule\noalign{}
\endhead
\bottomrule\noalign{}
\endlastfoot
无重力·探针 & 17.84 & 5.009 & 0.232 & 94.0\% & 27.2\% & 未全部通过 \\
带重力·无补偿（失败） & 1958.00 & 4.996 & 0.219 & 97.5\% & 22.7\% & 未全部通过 \\
带重力补偿·探针 & 23.31 & 5.009 & 0.208 & 96.5\% & 29.3\% & 未全部通过 \\
带重力补偿·夹爪 & 20.88 & 5.009 & 0.032 & 100.0\% & 95.4\% & 未全部通过 \\
\end{longtable}\endgroup

\FloatBarrier\begingroup\small\setlength{\tabcolsep}{2.5pt}
\begin{longtable}[]{@{}
  >{\raggedright\arraybackslash}p{(\linewidth - 4\tabcolsep) * \real{0.3249}}
  >{\raggedright\arraybackslash}p{(\linewidth - 4\tabcolsep) * \real{0.4296}}
  >{\raggedright\arraybackslash}p{(\linewidth - 4\tabcolsep) * \real{0.2456}}@{}}
\caption{重力及几何对照：设置与结果}\label{tab:s13-3}\tabularnewline
\toprule\noalign{}
\begin{minipage}[b]{\linewidth}\raggedright
条件
\end{minipage} & \begin{minipage}[b]{\linewidth}\raggedright
前三轴5s RMSE/度
\end{minipage} & \begin{minipage}[b]{\linewidth}\raggedright
首次周期接触/s
\end{minipage} \\
\midrule\noalign{}
\endfirsthead
\toprule\noalign{}
\begin{minipage}[b]{\linewidth}\raggedright
条件
\end{minipage} & \begin{minipage}[b]{\linewidth}\raggedright
前三轴5s RMSE/度
\end{minipage} & \begin{minipage}[b]{\linewidth}\raggedright
首次周期接触/s
\end{minipage} \\
\midrule\noalign{}
\endhead
\bottomrule\noalign{}
\endlastfoot
无重力·探针 & 0.0394 / 0.0266 / 0.0276 & 10.50 \\
带重力·无补偿（失败） & 0.0947 / 1.5800 / 0.1593 & 0.11 \\
带重力补偿·探针 & 0.0394 / 0.0266 / 0.0274 & 10.50 \\
带重力补偿·夹爪 & 0.0394 / 0.0266 / 0.0274 & 10.21 \\
\end{longtable}\endgroup

无补偿条件在初始轨迹阶段即下坠并碰到平面，属于控制失败，不能因为最后均力接近5N就判定可用。该条件的巨大峰值仅用于说明重力补偿缺失的后果，不纳入正常接触方案的改善百分比。补偿条件的前三轴误差保持在较小范围。

\stagefigure{figures/ch05/s13_gravity_f06ae4bc.png}{相同探针条件下，关闭重力与开启重力并补偿的对照。}{fig:s13-detail-1}


\paragraph{阶段结论}重力缺失补偿造成初态下坠；恢复手部与双指碰撞后重新建立对照，其他连杆碰撞仍未全面恢复。

\stagecode{s13}

\FloatBarrier

\section{接近减速、目标力缓升与修正速度}

\label{stage:s14}

\subsection{接近与力切换优化}\label{s14-detail-1}

\FloatBarrier\begingroup\small\setlength{\tabcolsep}{2.5pt}
\begin{longtable}[]{@{}
  >{\raggedright\arraybackslash}p{(\linewidth - 12\tabcolsep) * \real{0.1859}}
  >{\raggedright\arraybackslash}p{(\linewidth - 12\tabcolsep) * \real{0.1355}}
  >{\raggedright\arraybackslash}p{(\linewidth - 12\tabcolsep) * \real{0.1355}}
  >{\raggedright\arraybackslash}p{(\linewidth - 12\tabcolsep) * \real{0.1355}}
  >{\raggedright\arraybackslash}p{(\linewidth - 12\tabcolsep) * \real{0.1355}}
  >{\raggedright\arraybackslash}p{(\linewidth - 12\tabcolsep) * \real{0.1355}}
  >{\raggedright\arraybackslash}p{(\linewidth - 12\tabcolsep) * \real{0.1367}}@{}}
\caption{接近与力切换优化：设置与结果}\label{tab:s14-1}\tabularnewline
\toprule\noalign{}
\begin{minipage}[b]{\linewidth}\raggedright
条件
\end{minipage} & \begin{minipage}[b]{\linewidth}\raggedright
原始峰值/N
\end{minipage} & \begin{minipage}[b]{\linewidth}\raggedright
末段均力/N
\end{minipage} & \begin{minipage}[b]{\linewidth}\raggedright
力RMSE/N
\end{minipage} & \begin{minipage}[b]{\linewidth}\raggedright
周期带内率
\end{minipage} & \begin{minipage}[b]{\linewidth}\raggedright
原始带内率
\end{minipage} & \begin{minipage}[b]{\linewidth}\raggedright
四项检查
\end{minipage} \\
\midrule\noalign{}
\endfirsthead
\toprule\noalign{}
\begin{minipage}[b]{\linewidth}\raggedright
条件
\end{minipage} & \begin{minipage}[b]{\linewidth}\raggedright
原始峰值/N
\end{minipage} & \begin{minipage}[b]{\linewidth}\raggedright
末段均力/N
\end{minipage} & \begin{minipage}[b]{\linewidth}\raggedright
力RMSE/N
\end{minipage} & \begin{minipage}[b]{\linewidth}\raggedright
周期带内率
\end{minipage} & \begin{minipage}[b]{\linewidth}\raggedright
原始带内率
\end{minipage} & \begin{minipage}[b]{\linewidth}\raggedright
四项检查
\end{minipage} \\
\midrule\noalign{}
\endhead
\bottomrule\noalign{}
\endlastfoot
带重力补偿·夹爪 & 20.88 & 5.009 & 0.032 & 100.0\% & 95.4\% & 未全部通过 \\
夹爪＋目标力缓升 & 17.22 & 5.009 & 0.033 & 100.0\% & 95.8\% & 未全部通过 \\
接近减速初测 & 11.20 & 5.009 & 0.030 & 100.0\% & 97.0\% & 通过 \\
统一高度算法复测 & 21.25 & 5.009 & 0.034 & 100.0\% & 94.9\% & 未全部通过 \\
修正速度3mm/s & 16.87 & 5.009 & 0.027 & 100.0\% & 97.2\% & 未全部通过 \\
提前15mm减速 & 26.61 & 5.007 & 0.032 & 100.0\% & 95.9\% & 未全部通过 \\
减速＋修正速度3mm/s & 15.91 & 5.007 & 0.034 & 100.0\% & 95.2\% & 未全部通过 \\
附加法向阻尼100 & 16.85 & 5.009 & 0.031 & 100.0\% & 94.8\% & 未全部通过 \\
附加法向阻尼300 & 24.36 & 5.014 & 0.108 & 100.0\% & 59.2\% & 未全部通过 \\
阻尼100＋修正速度1mm/s & 25.17 & 5.007 & 0.027 & 100.0\% & 95.9\% & 未全部通过 \\
统一候选·修正3mm/s & 11.11 & 5.007 & 0.018 & 100.0\% & 100.0\% & 通过 \\
统一候选·修正30mm/s & 13.36 & 5.006 & 0.008 & 100.0\% & 100.0\% & 通过 \\
最终方案 & 11.11 & 5.007 & 0.018 & 100.0\% & 100.0\% & 通过 \\
\end{longtable}\endgroup

首次接触之前没有可用接触力，单靠接触门控无法预防碰撞。目标力缓升避免切换瞬间直接要求5N，提前减速则降低接触前速度，两者作用阶段不同。后续限制导纳修正速度和扩大减速范围，检查能否减少切换后的再次撞击。

初测gripper\_brake峰值较低，但名义平面高度计算改为不读取扰动值、增加运行检查后，默认条件峰值重新升高。两个条件并非逐位相同的复现，不能把初测较好数值当作可重复性能；保留两组结果以展示刚性接触对微小差异的敏感性。最终验证以统一代码下的新条件为准。

\stagefigure{figures/ch05/s14_approach_cb3b566a.png}{相同刚性夹爪场景下的接近与目标力切换对照。}{fig:s14-detail-1}

\stagefigure{figures/ch05/s14_refinement_0c4260a8.png}{修正速度与提前减速范围的后续对照。}{fig:s14-detail-2}


\paragraph{阶段结论}选用统一方案原始峰值11.11牛，末段均力约5.007牛；单因素和版本变化分别说明，不用较好先导替代复测。



% BEGIN SIM_VISUALS s14
\subsection{带重力接触的画面与阶段}
两侧均来自同一场景下保存的实际相机帧。相机画面能说明夹爪与平面的相对位置，不能独立分辨数牛的力差；因此保留同步力曲线，并在下面另外给出不含曲线的接近与接触画面。

\motionfigure{figures/ch05/panda_gravity_comparison_full_frame.png}{相同重力、几何与内环下的位置控制和优化导纳：实际画面与同步均力曲线。}{fig:motion-panda_gravity_comparison_full_frame}

\motionfigure{figures/ch05/panda_contact_camera_pair_frame000.png}{两种控制条件的纯相机画面，保存帧序号0（约0.001秒仿真时间）。}{fig:motion-panda_contact_camera_pair_frame000}

\motionfigure{figures/ch05/panda_contact_camera_pair_frame140.png}{两种控制条件的纯相机画面，保存帧序号140（约14.001秒仿真时间）。}{fig:motion-panda_contact_camera_pair_frame140}

\motionfigure{figures/ch05/panda_contact_camera_pair_frame300.png}{两种控制条件的纯相机画面，保存帧序号300（约30.001秒仿真时间）。}{fig:motion-panda_contact_camera_pair_frame300}
% END SIM_VISUALS s14
\stagecode{s14}
\stagevideo{videos/ch05/panda_gravity_contact_comparison.mp4}

\FloatBarrier

\section{扰动、接触材质与数值敏感性}

\label{stage:s15}

\subsection{环境与周期扰动}\label{s15-detail-1}

\FloatBarrier\begingroup\small\setlength{\tabcolsep}{2.5pt}
\begin{longtable}[]{@{}
  >{\raggedright\arraybackslash}p{(\linewidth - 12\tabcolsep) * \real{0.1861}}
  >{\raggedright\arraybackslash}p{(\linewidth - 12\tabcolsep) * \real{0.1356}}
  >{\raggedright\arraybackslash}p{(\linewidth - 12\tabcolsep) * \real{0.1356}}
  >{\raggedright\arraybackslash}p{(\linewidth - 12\tabcolsep) * \real{0.1356}}
  >{\raggedright\arraybackslash}p{(\linewidth - 12\tabcolsep) * \real{0.1356}}
  >{\raggedright\arraybackslash}p{(\linewidth - 12\tabcolsep) * \real{0.1356}}
  >{\raggedright\arraybackslash}p{(\linewidth - 12\tabcolsep) * \real{0.1356}}@{}}
\caption{环境与周期扰动：设置与结果}\label{tab:s15-1}\tabularnewline
\toprule\noalign{}
\begin{minipage}[b]{\linewidth}\raggedright
条件
\end{minipage} & \begin{minipage}[b]{\linewidth}\raggedright
原始峰值/N
\end{minipage} & \begin{minipage}[b]{\linewidth}\raggedright
末段均力/N
\end{minipage} & \begin{minipage}[b]{\linewidth}\raggedright
力RMSE/N
\end{minipage} & \begin{minipage}[b]{\linewidth}\raggedright
周期带内率
\end{minipage} & \begin{minipage}[b]{\linewidth}\raggedright
原始带内率
\end{minipage} & \begin{minipage}[b]{\linewidth}\raggedright
四项检查
\end{minipage} \\
\midrule\noalign{}
\endfirsthead
\toprule\noalign{}
\begin{minipage}[b]{\linewidth}\raggedright
条件
\end{minipage} & \begin{minipage}[b]{\linewidth}\raggedright
原始峰值/N
\end{minipage} & \begin{minipage}[b]{\linewidth}\raggedright
末段均力/N
\end{minipage} & \begin{minipage}[b]{\linewidth}\raggedright
力RMSE/N
\end{minipage} & \begin{minipage}[b]{\linewidth}\raggedright
周期带内率
\end{minipage} & \begin{minipage}[b]{\linewidth}\raggedright
原始带内率
\end{minipage} & \begin{minipage}[b]{\linewidth}\raggedright
四项检查
\end{minipage} \\
\midrule\noalign{}
\endhead
\bottomrule\noalign{}
\endlastfoot
统一候选·修正1mm/s & 16.13 & 5.008 & 0.010 & 100.0\% & 100.0\% & 未全部通过 \\
最终·摩擦1.0 & 8.86 & 5.007 & 0.008 & 100.0\% & 100.0\% & 通过 \\
最终·摩擦0.1 & 10.08 & 5.008 & 0.054 & 100.0\% & 45.8\% & 通过 \\
最终·物理步长0.5ms & 20.43 & 5.007 & 0.026 & 100.0\% & 93.4\% & 未全部通过 \\
最终·平面+1mm & 11.83 & 5.007 & 0.011 & 100.0\% & 100.0\% & 通过 \\
最终·平面+5mm & 12.83 & 5.004 & 0.006 & 100.0\% & 100.0\% & 通过 \\
最终·平面−1mm & 19.08 & 5.006 & 0.015 & 100.0\% & 100.0\% & 未全部通过 \\
最终·平面−5mm & 11.64 & 5.008 & 0.010 & 100.0\% & 100.0\% & 通过 \\
最终·重复 & 11.11 & 5.007 & 0.018 & 100.0\% & 100.0\% & 通过 \\
最终方案 & 11.11 & 5.007 & 0.018 & 100.0\% & 100.0\% & 通过 \\
最终·柔顺50kN/m & 5.16 & 5.007 & 0.013 & 100.0\% & 100.0\% & 通过 \\
柔顺·物理步长0.5ms & 5.11 & 5.007 & 0.014 & 100.0\% & 100.0\% & 通过 \\
柔顺·平面+5mm & 5.42 & 5.004 & 0.007 & 100.0\% & 100.0\% & 通过 \\
柔顺·平面−5mm & 5.09 & 5.008 & 0.010 & 100.0\% & 100.0\% & 通过 \\
20ms控制周期 & 21.78 & 5.009 & 0.032 & 100.0\% & 95.0\% & 未全部通过 \\
5ms控制周期 & 16.16 & 5.009 & 0.026 & 100.0\% & 97.2\% & 未全部通过 \\
默认方案·摩擦1.0 & 28.88 & 5.009 & 0.035 & 100.0\% & 94.5\% & 未全部通过 \\
默认方案·摩擦0.1 & 12.30 & 5.012 & 0.090 & 100.0\% & 38.6\% & 通过 \\
默认方案·平面+1mm & 27.43 & 5.009 & 0.029 & 100.0\% & 96.0\% & 未全部通过 \\
默认方案·平面+5mm & 24.75 & 5.009 & 0.033 & 100.0\% & 93.7\% & 未全部通过 \\
默认方案·平面−1mm & 14.64 & 5.009 & 0.035 & 100.0\% & 96.5\% & 通过 \\
默认方案·平面−5mm & 0.00 & 0.000 & 5.000 & 0.0\% & 0.0\% & 未全部通过 \\
质量矩阵补偿内环 & 62.49 & 5.017 & 0.205 & 97.5\% & 28.4\% & 未全部通过 \\
统一高度算法复测 & 21.25 & 5.009 & 0.034 & 100.0\% & 94.9\% & 未全部通过 \\
柔顺材质2kN/m & 5.32 & 5.007 & 0.011 & 100.0\% & 100.0\% & 通过 \\
柔顺材质50kN/m & 5.65 & 5.009 & 0.014 & 100.0\% & 100.0\% & 通过 \\
4次位置求解迭代 & 12.77 & 5.009 & 0.011 & 100.0\% & 100.0\% & 通过 \\
\end{longtable}\endgroup

高度扰动相对同一几何条件的名义平面施加，不给控制器额外告知实际偏差。摩擦变化只检验当前法向接触工况，不能证明切向擦拭或滑动稳定性。质量矩阵补偿使用加速度形式增益，与普通PID的单位和数值不同；这一组仅作为备选内环，不直接归因补偿方法优劣。5/20ms结果为周期敏感性测试，主实验仍以10ms为准。

\FloatBarrier\begingroup\small\setlength{\tabcolsep}{2.5pt}
\begin{longtable}[]{@{}
  >{\raggedright\arraybackslash}p{(\linewidth - 12\tabcolsep) * \real{0.1680}}
  >{\raggedright\arraybackslash}p{(\linewidth - 12\tabcolsep) * \real{0.1387}}
  >{\raggedright\arraybackslash}p{(\linewidth - 12\tabcolsep) * \real{0.1387}}
  >{\raggedright\arraybackslash}p{(\linewidth - 12\tabcolsep) * \real{0.1387}}
  >{\raggedright\arraybackslash}p{(\linewidth - 12\tabcolsep) * \real{0.1387}}
  >{\raggedright\arraybackslash}p{(\linewidth - 12\tabcolsep) * \real{0.1387}}
  >{\raggedright\arraybackslash}p{(\linewidth - 12\tabcolsep) * \real{0.1387}}@{}}
\caption{环境与周期扰动：设置与结果}\label{tab:s15-2}\tabularnewline
\toprule\noalign{}
\begin{minipage}[b]{\linewidth}\raggedright
条件
\end{minipage} & \begin{minipage}[b]{\linewidth}\raggedright
原始峰值/N
\end{minipage} & \begin{minipage}[b]{\linewidth}\raggedright
末段均力/N
\end{minipage} & \begin{minipage}[b]{\linewidth}\raggedright
力RMSE/N
\end{minipage} & \begin{minipage}[b]{\linewidth}\raggedright
周期带内率
\end{minipage} & \begin{minipage}[b]{\linewidth}\raggedright
原始带内率
\end{minipage} & \begin{minipage}[b]{\linewidth}\raggedright
四项检查
\end{minipage} \\
\midrule\noalign{}
\endfirsthead
\toprule\noalign{}
\begin{minipage}[b]{\linewidth}\raggedright
条件
\end{minipage} & \begin{minipage}[b]{\linewidth}\raggedright
原始峰值/N
\end{minipage} & \begin{minipage}[b]{\linewidth}\raggedright
末段均力/N
\end{minipage} & \begin{minipage}[b]{\linewidth}\raggedright
力RMSE/N
\end{minipage} & \begin{minipage}[b]{\linewidth}\raggedright
周期带内率
\end{minipage} & \begin{minipage}[b]{\linewidth}\raggedright
原始带内率
\end{minipage} & \begin{minipage}[b]{\linewidth}\raggedright
四项检查
\end{minipage} \\
\midrule\noalign{}
\endhead
\bottomrule\noalign{}
\endlastfoot
接近减速初测 & 11.20 & 5.009 & 0.030 & 100.0\% & 97.0\% & 通过 \\
柔顺材质10kN/m & 6.40 & 5.009 & 0.013 & 100.0\% & 100.0\% & 通过 \\
柔顺材质2kN/m & 5.32 & 5.007 & 0.011 & 100.0\% & 100.0\% & 通过 \\
柔顺材质50kN/m & 5.65 & 5.009 & 0.014 & 100.0\% & 100.0\% & 通过 \\
\end{longtable}\endgroup

柔顺材质的峰值与末段波动较小，说明接触物理模型会明显影响结果。软化环境不等于刚性控制问题已解决，最终刚性与柔顺条件分别给出。

\stagefigure{figures/ch05/s15_material_4e8027ac.png}{刚性与柔顺材质对照；此处改变环境，不计算跨材质算法改善比例。}{fig:s15-detail-1}

\stagefigure{figures/ch05/s15_validation_7338073d.png}{默认减速方案的高度与摩擦扰动。}{fig:s15-detail-2}

\stagefigure{figures/ch05/s15_timestep_92c24626.png}{物理步长减半的数值检查；原始峰值与周期均力分别解释。}{fig:s15-detail-3}

\stagefigure{figures/ch05/s15_final_validation_ebdf4225.png}{统一修正方案的独立条件验证。}{fig:s15-detail-4}

\FloatBarrier\begingroup\small\setlength{\tabcolsep}{2.5pt}
\begin{longtable}[]{@{}
  >{\raggedright\arraybackslash}p{(\linewidth - 12\tabcolsep) * \real{0.1728}}
  >{\raggedright\arraybackslash}p{(\linewidth - 12\tabcolsep) * \real{0.1379}}
  >{\raggedright\arraybackslash}p{(\linewidth - 12\tabcolsep) * \real{0.1379}}
  >{\raggedright\arraybackslash}p{(\linewidth - 12\tabcolsep) * \real{0.1379}}
  >{\raggedright\arraybackslash}p{(\linewidth - 12\tabcolsep) * \real{0.1379}}
  >{\raggedright\arraybackslash}p{(\linewidth - 12\tabcolsep) * \real{0.1379}}
  >{\raggedright\arraybackslash}p{(\linewidth - 12\tabcolsep) * \real{0.1379}}@{}}
\caption{环境与周期扰动：设置与结果}\label{tab:s15-3}\tabularnewline
\toprule\noalign{}
\begin{minipage}[b]{\linewidth}\raggedright
条件
\end{minipage} & \begin{minipage}[b]{\linewidth}\raggedright
原始峰值/N
\end{minipage} & \begin{minipage}[b]{\linewidth}\raggedright
末段均力/N
\end{minipage} & \begin{minipage}[b]{\linewidth}\raggedright
力RMSE/N
\end{minipage} & \begin{minipage}[b]{\linewidth}\raggedright
周期带内率
\end{minipage} & \begin{minipage}[b]{\linewidth}\raggedright
原始带内率
\end{minipage} & \begin{minipage}[b]{\linewidth}\raggedright
四项检查
\end{minipage} \\
\midrule\noalign{}
\endfirsthead
\toprule\noalign{}
\begin{minipage}[b]{\linewidth}\raggedright
条件
\end{minipage} & \begin{minipage}[b]{\linewidth}\raggedright
原始峰值/N
\end{minipage} & \begin{minipage}[b]{\linewidth}\raggedright
末段均力/N
\end{minipage} & \begin{minipage}[b]{\linewidth}\raggedright
力RMSE/N
\end{minipage} & \begin{minipage}[b]{\linewidth}\raggedright
周期带内率
\end{minipage} & \begin{minipage}[b]{\linewidth}\raggedright
原始带内率
\end{minipage} & \begin{minipage}[b]{\linewidth}\raggedright
四项检查
\end{minipage} \\
\midrule\noalign{}
\endhead
\bottomrule\noalign{}
\endlastfoot
最终·柔顺50kN/m & 5.16 & 5.007 & 0.013 & 100.0\% & 100.0\% & 通过 \\
柔顺·平面+5mm & 5.42 & 5.004 & 0.007 & 100.0\% & 100.0\% & 通过 \\
柔顺·平面−5mm & 5.09 & 5.008 & 0.010 & 100.0\% & 100.0\% & 通过 \\
\end{longtable}\endgroup

在最终控制器保持不变的情况下另测50kN/m、150N·s/m的等效柔顺接触，表示可能通过末端薄缓冲层改变接触物理。该结果需要真实材料与接触面积标定后才能迁移，不能作为刚性环境问题已经解决的依据。

\subsubsection{物理步长与冲量检查}\label{s15-detail-2}

\FloatBarrier\begingroup\small\setlength{\tabcolsep}{2.5pt}
\begin{longtable}[]{@{}
  >{\raggedright\arraybackslash}p{(\linewidth - 8\tabcolsep) * \real{0.2409}}
  >{\raggedright\arraybackslash}p{(\linewidth - 8\tabcolsep) * \real{0.1853}}
  >{\raggedright\arraybackslash}p{(\linewidth - 8\tabcolsep) * \real{0.1853}}
  >{\raggedright\arraybackslash}p{(\linewidth - 8\tabcolsep) * \real{0.2032}}
  >{\raggedright\arraybackslash}p{(\linewidth - 8\tabcolsep) * \real{0.1853}}@{}}
\caption{物理步长与冲量检查：设置与结果}\label{tab:s15-4}\tabularnewline
\toprule\noalign{}
\begin{minipage}[b]{\linewidth}\raggedright
条件
\end{minipage} & \begin{minipage}[b]{\linewidth}\raggedright
物理步长/ms
\end{minipage} & \begin{minipage}[b]{\linewidth}\raggedright
原始峰值/N
\end{minipage} & \begin{minipage}[b]{\linewidth}\raggedright
最大单步冲量/N·s
\end{minipage} & \begin{minipage}[b]{\linewidth}\raggedright
10ms均力峰值/N
\end{minipage} \\
\midrule\noalign{}
\endfirsthead
\toprule\noalign{}
\begin{minipage}[b]{\linewidth}\raggedright
条件
\end{minipage} & \begin{minipage}[b]{\linewidth}\raggedright
物理步长/ms
\end{minipage} & \begin{minipage}[b]{\linewidth}\raggedright
原始峰值/N
\end{minipage} & \begin{minipage}[b]{\linewidth}\raggedright
最大单步冲量/N·s
\end{minipage} & \begin{minipage}[b]{\linewidth}\raggedright
10ms均力峰值/N
\end{minipage} \\
\midrule\noalign{}
\endhead
\bottomrule\noalign{}
\endlastfoot
最终方案 & 1.0 & 11.11 & 0.01111 & 5.143 \\
最终·物理步长0.5ms & 0.5 & 20.43 & 0.01022 & 5.085 \\
最终·柔顺50kN/m & 1.0 & 5.16 & 0.00516 & 5.147 \\
柔顺·物理步长0.5ms & 0.5 & 5.11 & 0.00255 & 5.084 \\
\end{longtable}\endgroup

刚性条件把物理步长从1ms减到0.5ms后，原始峰值从11.11N变成20.43N，末段原始带内率也下降，说明瞬时峰值尚未数值收敛。单步接触力来自求解器冲量除以步长，因而还需比较冲量和固定分箱均力，不能把离散尖峰直接解释为真实传感器的力上界。改变步长也会改变启动预步进的微小初始误差，这一组属于数值敏感性检查，不是严格的连续系统收敛证明。

有限刚度条件下，1ms与0.5ms的原始峰值分别为5.16N、5.11N，说明在该阶段范围内对步长更不敏感。真实迁移仍需要标定接触刚度、阻尼、接触面积和测力带宽。


\subsection{结果边界和失败分析}\label{s15-detail-3}

\FloatBarrier\begingroup\small\setlength{\tabcolsep}{2.5pt}
\begin{longtable}[]{@{}
  >{\raggedright\arraybackslash}p{(\linewidth - 2\tabcolsep) * \real{0.5676}}
  >{\raggedright\arraybackslash}p{(\linewidth - 2\tabcolsep) * \real{0.4324}}@{}}
\caption{结果边界和失败分析：设置与结果}\label{tab:s15-5}\tabularnewline
\toprule\noalign{}
\begin{minipage}[b]{\linewidth}\raggedright
选用条件的指标
\end{minipage} & \begin{minipage}[b]{\linewidth}\raggedright
结果
\end{minipage} \\
\midrule\noalign{}
\endfirsthead
\toprule\noalign{}
\begin{minipage}[b]{\linewidth}\raggedright
选用条件的指标
\end{minipage} & \begin{minipage}[b]{\linewidth}\raggedright
结果
\end{minipage} \\
\midrule\noalign{}
\endhead
\bottomrule\noalign{}
\endlastfoot
条件 & final\_selected \\
1ms原始力峰值 & 11.112 N \\
末2s周期均力带内率 & 100.0\% \\
末2s原始力带内率 & 100.0\% \\
末2s原始峰值 & 5.255 N \\
最长连续周期带内跨度 & 14.53 s \\
\end{longtable}\endgroup

该阶段工程检查设为：原始峰值\(\leq\)15N、末2s周期带内率\(\geq\)99\%、均力RMSE\(\leq\)0.25N、前三轴各自5s RMSE\(\leq\)0.1°。这些是本项目用于比较的目标，不是工业标准或设备认证。表中的通过表示四项同时满足，完成标记只表示仿真正常结束。

接触后的短促原始力脉冲仍可能超过5±0.5N。平面高度、摩擦、求解器和接触几何变化都可能放大这些脉冲；没有将刚性求解器的接触脉冲等同于真实六维力传感器输出。尤其低摩擦条件下，即使周期均力合格，原始力仍可能频繁越界。更长时间、多姿态、多随机扰动及远程网络和实体试验仍需要补充。

新增记录区分了原始合力、周期均力、反馈EMA和目标力；因此可以检查控制器是否因为滤波延迟而反复接触。重复运行用于检查确定性和实现差异，不用于宣称统计显著性。


\paragraph{阶段结论}刚性接触减半步长后峰值由11.11增至20.43牛，瞬时峰值未数值收敛；柔顺材质另列，不作跨材质算法提升比较。

\stagecode{s15}

\FloatBarrier

\section{动作块执行、延迟与跨进程通信}

\label{stage:s16}

\subsection{扩展：学习策略动作执行接口}\label{s16-detail-1}

扩展目标是检查策略返回动作以后，执行端如何处理动作时标、突变和断流。参考LeRobot预测与执行解耦、动作队列的设计；该阶段为独立的确定性接口测试，没有接入ACT、SmolVLA或WAM推理，也没有实现训练式RTC。

\href{https://huggingface.co/docs/lerobot/async}{参考：LeRobot异步推理与动作队列}

七组测试均在Isaac Panda实际动力学中执行，普通PID内环不变。前三轴使用三倍幅度的5秒五次参考，固定面下移以排除接触。模拟策略每200ms生成40步未来动作，动作间隔10ms，固定送达延迟120ms。这样可以把执行机制的影响与模型预测错误分开。

\FloatBarrier\begingroup\small\setlength{\tabcolsep}{2.5pt}
\begin{longtable}[]{@{}
  >{\raggedright\arraybackslash}p{(\linewidth - 2\tabcolsep) * \real{0.3529}}
  >{\raggedright\arraybackslash}p{(\linewidth - 2\tabcolsep) * \real{0.6471}}@{}}
\caption{扩展：学习策略动作执行接口：设置与结果}\label{tab:s16-1}\tabularnewline
\toprule\noalign{}
\begin{minipage}[b]{\linewidth}\raggedright
条件
\end{minipage} & \begin{minipage}[b]{\linewidth}\raggedright
注入或执行方式
\end{minipage} \\
\midrule\noalign{}
\endfirsthead
\toprule\noalign{}
\begin{minipage}[b]{\linewidth}\raggedright
条件
\end{minipage} & \begin{minipage}[b]{\linewidth}\raggedright
注入或执行方式
\end{minipage} \\
\midrule\noalign{}
\endhead
\bottomrule\noalign{}
\endlastfoot
真值定位 & 100Hz直接发送理想参考 \\
delayed\_naive & 动作块到达后从第0步执行，忽略生成时间 \\
delayed\_timestamp & 按当前时间减生成时间选择动作，跳过过期前缀 \\
jitter\_raw & 每个块在第一轴交替叠加±0.04rad偏差 \\
jitter\_limited & 同样偏差，命令速度限制0.35rad/s、加速度1rad/s² \\
outage\_stale & 2---3秒不生成新块，继续消费旧块并停留在最后目标 \\
outage\_hold & 同样断流，源数据年龄超过0.5秒后保持一次读取的实测姿态，恢复时仍限速限加速度 \\
\end{longtable}\endgroup

\FloatBarrier\begingroup\small\setlength{\tabcolsep}{2.5pt}
\begin{longtable}[]{@{}
  >{\raggedright\arraybackslash}p{(\linewidth - 6\tabcolsep) * \real{0.2934}}
  >{\raggedright\arraybackslash}p{(\linewidth - 6\tabcolsep) * \real{0.2355}}
  >{\raggedright\arraybackslash}p{(\linewidth - 6\tabcolsep) * \real{0.2355}}
  >{\raggedright\arraybackslash}p{(\linewidth - 6\tabcolsep) * \real{0.2355}}@{}}
\caption{扩展：学习策略动作执行接口：设置与结果}\label{tab:s16-2}\tabularnewline
\toprule\noalign{}
\begin{minipage}[b]{\linewidth}\raggedright
执行条件
\end{minipage} & \begin{minipage}[b]{\linewidth}\raggedright
实测对理想RMSE/度
\end{minipage} & \begin{minipage}[b]{\linewidth}\raggedright
最大命令跳变/度
\end{minipage} & \begin{minipage}[b]{\linewidth}\raggedright
峰值力矩/Nm
\end{minipage} \\
\midrule\noalign{}
\endfirsthead
\toprule\noalign{}
\begin{minipage}[b]{\linewidth}\raggedright
执行条件
\end{minipage} & \begin{minipage}[b]{\linewidth}\raggedright
实测对理想RMSE/度
\end{minipage} & \begin{minipage}[b]{\linewidth}\raggedright
最大命令跳变/度
\end{minipage} & \begin{minipage}[b]{\linewidth}\raggedright
峰值力矩/Nm
\end{minipage} \\
\midrule\noalign{}
\endhead
\bottomrule\noalign{}
\endlastfoot
真值定位 & 0.095 & 0.039 & 0.062 \\
delayed\_naive & 0.330 & 0.039 & 0.062 \\
delayed\_timestamp & 0.095 & 0.039 & 0.062 \\
jitter\_raw & 1.145 & 4.622 & 17.795 \\
jitter\_limited & 0.946 & 0.146 & 0.879 \\
outage\_stale & 0.532 & 2.726 & 9.510 \\
outage\_hold & 1.031 & 0.201 & 3.052 \\
\end{longtable}\endgroup

RMSE统计前三轴相对理想参考的误差；跳变及峰值力矩统计全部七轴。所有条件使用同一条已知轨迹，没有多随机种子统计。命令限幅经过逐样本检查，但这不等于实际关节速度限幅或碰撞安全证明。

\stagefigure{figures/ch05/s16_execution_delay_fa537d74.png}{固定120ms延迟下，忽略时标与按时标选取动作的比较。}{fig:s16-detail-1}

忽略时标时，实际轨迹RMSE为0.330°；按时标取动作后为0.095°。这里的未来动作来自已知解析参考，时标对齐能够恢复理想执行；真实VLA预测误差、观察过时和环境变化并不能靠跳过前缀消除。

\stagefigure{figures/ch05/s16_execution_jitter_307c270b.png}{同样的动作块偏差下，直接执行与命令限速限加速度。}{fig:s16-detail-2}

限幅将最大命令跳变从4.622°降至0.146°，峰值力矩从17.795Nm降至0.879Nm。平滑会改变原动作时序，因此不能只看动作更顺，还应核对跟踪误差及任务阶段。

\stagefigure{figures/ch05/s16_execution_outage_aac306ae.png}{橙色区间为停止产生动作块的2—3秒；包含送达延迟与超时等待。}{fig:s16-detail-3}

超时保持与受限恢复使最大命令跳变从2.726°降至0.201°，但轨迹RMSE从0.532°增加到1.031°。这是有意暂停追赶带来的代价，不能写成所有指标都改善。超时触发1次；该保持机制尚未在重力、接触负载或真机上验证。

\subsubsection{动作窗口覆盖的接口回放}\label{s16-detail-2}

另外进行54组接口级回放：延迟40/120/240ms、块更新周期100/200/400ms、动作窗口200/400/800ms，分别比较忽略时标和按时标消费。这里只计算动作命令与窗口覆盖，不再次运行机器人动力学，不计入物理仿真次数。

\FloatBarrier\begingroup\small\setlength{\tabcolsep}{2.5pt}
\begin{longtable}[]{@{}
  >{\raggedright\arraybackslash}p{(\linewidth - 6\tabcolsep) * \real{0.2500}}
  >{\raggedright\arraybackslash}p{(\linewidth - 6\tabcolsep) * \real{0.2500}}
  >{\raggedright\arraybackslash}p{(\linewidth - 6\tabcolsep) * \real{0.2500}}
  >{\raggedright\arraybackslash}p{(\linewidth - 6\tabcolsep) * \real{0.2500}}@{}}
\caption{动作窗口覆盖的接口回放：设置与结果}\label{tab:s16-3}\tabularnewline
\toprule\noalign{}
\begin{minipage}[b]{\linewidth}\raggedright
延迟/ms
\end{minipage} & \begin{minipage}[b]{\linewidth}\raggedright
更新周期/ms
\end{minipage} & \begin{minipage}[b]{\linewidth}\raggedright
窗口/ms
\end{minipage} & \begin{minipage}[b]{\linewidth}\raggedright
窗口过期时间占比
\end{minipage} \\
\midrule\noalign{}
\endfirsthead
\toprule\noalign{}
\begin{minipage}[b]{\linewidth}\raggedright
延迟/ms
\end{minipage} & \begin{minipage}[b]{\linewidth}\raggedright
更新周期/ms
\end{minipage} & \begin{minipage}[b]{\linewidth}\raggedright
窗口/ms
\end{minipage} & \begin{minipage}[b]{\linewidth}\raggedright
窗口过期时间占比
\end{minipage} \\
\midrule\noalign{}
\endhead
\bottomrule\noalign{}
\endlastfoot
40 & 200 & 400 & 0.0\% \\
120 & 200 & 200 & 59.0\% \\
120 & 200 & 400 & 0.0\% \\
240 & 200 & 400 & 19.3\% \\
\end{longtable}\endgroup

\stagefigure{figures/ch05/s16_stream_coverage_1c419e63.png}{动作窗口不足时，即使按时标消费也会用完可执行动作。}{fig:s16-detail-4}

在固定延迟、无丢包的本模型下，窗口至少覆盖送达延迟与下一次块更新的间隔，才能持续提供未过期动作。真实系统还需要为抖动和计算时间变化留余量。扩大窗口可能降低反应速度，不能无限增加。配置化后的执行器还逐帧重放了七组已执行命令，与原命令一致。

相关代码在 \path|scripts/action_stream.py|、\path|scripts/analyze_execution.py| 和 \path|scripts/sweep_stream.py|。结果分别在 \path|execution_results.json|、\path|stream_sweep.json|。下一步接入实际策略的带时标动作块，在相同扰动下复测；该阶段不据此声称VLA已经能够完成操作任务。


\subsection{本机UDP跨进程通信}\label{s16-detail-3}

\FloatBarrier\begingroup\small\setlength{\tabcolsep}{2.5pt}
\begin{longtable}[]{@{}
  >{\raggedright\arraybackslash}p{(\linewidth - 8\tabcolsep) * \real{0.2037}}
  >{\raggedright\arraybackslash}p{(\linewidth - 8\tabcolsep) * \real{0.1959}}
  >{\raggedright\arraybackslash}p{(\linewidth - 8\tabcolsep) * \real{0.2052}}
  >{\raggedright\arraybackslash}p{(\linewidth - 8\tabcolsep) * \real{0.1993}}
  >{\raggedright\arraybackslash}p{(\linewidth - 8\tabcolsep) * \real{0.1959}}@{}}
\caption{本机UDP跨进程通信：设置与结果}\label{tab:s16-4}\tabularnewline
\toprule\noalign{}
\begin{minipage}[b]{\linewidth}\raggedright
工况
\end{minipage} & \begin{minipage}[b]{\linewidth}\raggedright
收到包数
\end{minipage} & \begin{minipage}[b]{\linewidth}\raggedright
收包轮询延迟P95/ms
\end{minipage} & \begin{minipage}[b]{\linewidth}\raggedright
生成到读取P95/ms
\end{minipage} & \begin{minipage}[b]{\linewidth}\raggedright
保持采样数
\end{minipage} \\
\midrule\noalign{}
\endfirsthead
\toprule\noalign{}
\begin{minipage}[b]{\linewidth}\raggedright
工况
\end{minipage} & \begin{minipage}[b]{\linewidth}\raggedright
收到包数
\end{minipage} & \begin{minipage}[b]{\linewidth}\raggedright
收包轮询延迟P95/ms
\end{minipage} & \begin{minipage}[b]{\linewidth}\raggedright
生成到读取P95/ms
\end{minipage} & \begin{minipage}[b]{\linewidth}\raggedright
保持采样数
\end{minipage} \\
\midrule\noalign{}
\endhead
\bottomrule\noalign{}
\endlastfoot
正常 & 80 & 10.07 & 10.21 & 1 \\
发送端延迟40ms & 79 & 10.04 & 50.23 & 5 \\
发送端停发1s & 60 & 10.10 & 10.24 & 57 \\
\end{longtable}\endgroup

发送与接收运行在独立进程，实际通过127.0.0.1 UDP传递未来动作块；发送20Hz、接收目标100Hz。延迟40ms和停发1s均由发送应用注入，动作来自解析五次轨迹，不运行VLA，也不驱动物理仿真或实体机器人。接收端按时间戳取动作，数据年龄超过0.5s时保持输出。

收包延迟包含100Hz轮询和操作系统调度，不能当作纯网络传输耗时或RTT。同机共享单调时钟；跨主机时钟同步、重排、丢包恢复和硬实时尚未验证。测试在该阶段仿真负载下完成，仅证明独立进程之间的实际通信和超时逻辑。相关代码在 scripts/network.py。


\paragraph{阶段结论}动作源是解析参考，验证时间戳、限幅、超时与恢复；这些结果不能解释为VLA网络控制成功。

\stagecode{s16}

\FloatBarrier
